\documentclass[12pt,a4paper]{article}

\usepackage{amsmath,amssymb,amsthm}
\usepackage{geometry}
\usepackage{hyperref}
\usepackage{booktabs}
\usepackage{graphicx}
\usepackage{xcolor}
\usepackage{cite}

\geometry{margin=2.5cm}

\newtheorem{theorem}{Theorem}
\newtheorem{lemma}[theorem]{Lemma}
\newtheorem{proposition}[theorem]{Proposition}
\newtheorem{definition}{Definition}
\newtheorem{remark}{Remark}

\title{\textbf{Differential Nonlinear Robustness of Critical States\\
in Fibonacci and Tribonacci Substitution Chains}\\[0.5em]
\large A Numerical Study of Discrete Nonlinear Schr\"{o}dinger Dynamics}

\author{Pablo Nogueira Grossi\\
G6 LLC, Newark, New Jersey, USA\\
ORCID: 0009-0000-6496-2186\\
\texttt{pablogrossi@hotmail.com}\\[0.3em]
\small Zenodo: \href{https://doi.org/10.5281/zenodo.20026943}{10.5281/zenodo.20026943}}

\date{May 2026}

\begin{document}

\maketitle

\begin{abstract}
We study the discrete nonlinear Schr\"{o}dinger (DNLS) equation on two
quasiperiodic tight-binding chains: the Fibonacci chain ($n=2$) and the
Rauzy--tribonacci chain ($n=3$), generated by their respective substitution
rules. Starting from mid-gap eigenstates of the linear Hamiltonian, we
integrate the DNLS time evolution over a range of nonlinearity strengths
$\lambda \in [0,10]$ and track the inverse participation ratio (IPR) as a
measure of localization. The tribonacci chain exhibits a mid-gap IPR
approximately four times larger than the Fibonacci chain in the linear limit,
consistent with the stronger multifractality of Rauzy fractal eigenstates.
Under nonlinear perturbation, the Fibonacci mid-gap state delocalizes rapidly
(IPR drops $\sim 60\%$ at $\lambda = 1.5$), while the tribonacci state
remains nearly pinned (IPR drop $< 5\%$ at the same coupling). We term this
\emph{differential nonlinear robustness}: the tribonacci chain's more
strongly multifractal spectrum provides greater effective resistance to
nonlinear delocalization. The characteristic decay constant of the
tribonacci amplitude envelope, $\eta \approx 1.839287$, is the
Perron--Frobenius eigenvalue of the tribonacci companion matrix, whose
existence and key properties ($\eta > 1$, strict antitonicity of $\eta^{-k}$)
are formally verified in Lean~4/Mathlib4 in the companion AXLE repository
\cite{AXLE2026}. Our results constitute, to our knowledge, the first
numerical study of DNLS dynamics on a tribonacci substitution chain, and
identify a qualitative difference in nonlinear response between $n=2$ and
$n=3$ substitution-generated quasicrystal models.
\end{abstract}

\textbf{Keywords:} tribonacci chain; Fibonacci quasicrystal; discrete
nonlinear Schr\"{o}dinger equation; multifractality; inverse participation
ratio; Rauzy substitution; nonlinear localization; quasiperiodic tight-binding

\section{Introduction}

Quasiperiodic tight-binding chains generated by substitution rules have been
a canonical setting for studying the interplay between long-range order,
multifractality, and localization since the 1980s. The Fibonacci chain---the
$n=2$ case of the $n$-bonacci family---occupies a central role: its spectrum
is a Cantor set of measure zero, its eigenstates are critical (neither
extended nor localized), and their multifractal dimensions are precisely
computable \cite{Kohmoto1983,Ostlund1983,Luck1989,Mace2017}. Topological
phases, charge pumping, and one-dimensional quasicrystalline superconductivity
have been analyzed in the Fibonacci setting \cite{Kraus2012,Lang2012}.

The natural generalization to tribonacci and higher $n$-bonacci substitution
chains is more recent. Krebbekx and collaborators \cite{Krebbekx2023}
introduced two tribonacci tight-binding models (hopping and on-site
modulation) via Rauzy substitutions, established that their spectra inherit
Rauzy-fractal internal space structure, and computed multifractal dimensions
showing stronger spectral clustering than in the Fibonacci case.
Higher-$n$ generalizations appear in the context of Rauzy tilings and
cut-and-project schemes, but exclusively in the linear (Schr\"{o}dinger
operator) setting \cite{Bellissard1989,Damanik2008}.

The nonlinear regime is comparatively unexplored. For the Fibonacci chain,
discrete nonlinear Schr\"{o}dinger (DNLS) dynamics have been studied in
photonic quasicrystals \cite{Lahini2009}, edge-soliton pumping has been
demonstrated \cite{Jurgensen2021}, and phason dynamics under nonlinear
perturbation have been analyzed \cite{Verbin2015}. For $n \geq 3$
substitution chains, however, no published work has placed an $n$-bonacci
recurrence structure inside a DNLS or Gross--Pitaevskii nonlinearity, nor
analyzed the resulting localization dynamics. This gap motivates the present
study.

Our central question is: \emph{does the stronger multifractality of the
tribonacci chain translate into qualitatively different nonlinear localization
behavior?} We answer affirmatively through direct numerical integration,
finding that tribonacci mid-gap states are substantially more robust against
nonlinear delocalization than their Fibonacci counterparts.

A secondary motivation comes from the dm$^3$ contact-geometric framework
\cite{Grossi2026_Vol1,Grossi2026_Wavenumber6}, where the tribonacci
recurrence $w(k+3) = w(k+2) + w(k+1) + w(k)$ arises from the three-axis
decomposition of a contact 3-manifold, and the geometric weight $\eta^{-k}$
with $\eta \approx 1.839287$ (the tribonacci constant) is the natural scaling
factor. The formal verification that $\eta > 1$ and that the weight sequence
$\{\eta^{-k}\}$ is strictly antitone---establishing that the amplitude
envelope is well-defined and decaying---is carried out in Lean~4/Mathlib4 in
the AXLE repository \cite{AXLE2026}, specifically in
\texttt{TribonacciMeasure.lean}. We cite this formal result in Section~3 as
the machine-verified foundation for the decay envelope used in our ansatz.

\section{The $n$-Bonacci Recurrence and Substitution Chains}

\subsection{The $n$-Bonacci Recurrence}

\begin{definition}[$n$-Bonacci Sequence]
For $n \geq 2$, the $n$-bonacci sequence $\{A_k^{(n)}\}_{k \geq 0}$ is
defined by
\[
  A_k^{(n)} = \sum_{i=1}^{n} A_{k-i}^{(n)}, \quad k \geq n,
\]
with initial conditions $A_0^{(n)} = \cdots = A_{n-2}^{(n)} = 0$,
$A_{n-1}^{(n)} = 1$.
\end{definition}

The characteristic polynomial is
\[
  p_n(x) = x^n - x^{n-1} - x^{n-2} - \cdots - x - 1 = 0.
\]
By the Perron--Frobenius theorem applied to the companion matrix
\[
C_n = \begin{pmatrix}
0 & 1 & 0 & \cdots & 0 \\
0 & 0 & 1 & \cdots & 0 \\
\vdots & & & \ddots & \vdots \\
0 & 0 & 0 & \cdots & 1 \\
1 & 1 & 1 & \cdots & 1
\end{pmatrix} \in M_n(\mathbb{Z}),
\]
there exists a unique dominant real root $\rho_n > 1$, which we call the
\emph{$n$-bonacci constant}. The sequence $\rho_n$ is strictly increasing
and satisfies $\rho_n \to 2$ monotonically. Numerical values:

\begin{center}
\begin{tabular}{ccl}
\toprule
$n$ & $\rho_n$ (6 d.p.) & Name \\
\midrule
2 & 1.618034 & Golden ratio $\varphi$ \\
3 & 1.839287 & Tribonacci constant $\eta$ \\
4 & 1.927562 & Tetranacci constant \\
5 & 1.965948 & Pentanacci constant \\
\bottomrule
\end{tabular}
\end{center}

\begin{remark}[Formal Verification]
For $n = 3$: the existence of $\eta \approx 1.839287$ as the unique real root
of $x^3 - x^2 - x - 1 = 0$ in $[1,2]$, together with $\eta > 1$ and the
strict antitonicity of the weight sequence $\{\eta^{-k}\}_{k \geq 0}$, is
formally proved in Lean~4/Mathlib4 without \texttt{sorry} in
\texttt{TribonacciMeasure.lean} of the AXLE repository \cite{AXLE2026}. The
proof uses the intermediate value theorem for the characteristic polynomial
(IVT witnesses: $p_3(1) = -2 < 0$, $p_3(2) = 1 > 0$) and the fact that
$\eta^{-1} \in (0,1)$ implies $\{\eta^{-k}\}$ is strictly decreasing. This
constitutes machine-verified backing for the decay rate used in the amplitude
ansatz (Section~3.2).
\end{remark}

\subsection{Substitution Chains}

The Fibonacci chain ($n=2$) is generated by the substitution $\sigma_2$:
\[
  \sigma_2: A \mapsto AB,\quad B \mapsto A.
\]
Starting from $A$, repeated application generates the one-sided Fibonacci
word $ABAABABAABAAB\cdots$, a sequence over the two-letter alphabet $\{A, B\}$.

The tribonacci (Rauzy) chain ($n=3$) is generated by the substitution
$\sigma_3$ \cite{Krebbekx2023,Rauzy1982}:
\[
  \sigma_3: 1 \mapsto 12,\quad 2 \mapsto 13,\quad 3 \mapsto 1.
\]
Starting from $1$, repeated application generates the tribonacci word over
$\{1, 2, 3\}$. The letter frequencies converge to the left Perron--Frobenius
eigenvector of the substitution matrix
\[
  M_3 = \begin{pmatrix} 1 & 1 & 1 \\ 1 & 0 & 0 \\ 0 & 1 & 0 \end{pmatrix},
\]
giving asymptotic frequencies $(f_1, f_2, f_3) \approx (0.5437, 0.2956, 0.1607)$.
We verified this numerically on a 600-symbol word (observed: $0.5433$,
$0.2967$, $0.1600$), confirming correct substitution implementation.

\subsection{Tight-Binding Hamiltonian}

We assign hopping amplitudes to letters via a modulation parameter $t \in
(0,1)$:
\[
  t_A = 1.0,\quad t_B = t,\quad t_C = t^2.
\]
We use $t = 0.5$ throughout, a value in the generic incommensurate regime
studied in \cite{Krebbekx2023}. The Hamiltonian on a chain of $N$ sites is
the tridiagonal matrix
\[
  H_{j,j+1} = H_{j+1,j} = t_{\sigma(j)}, \quad H_{jj} = 0,
\]
where $\sigma(j)$ is the $j$-th letter of the substitution word. For $N = 500$
sites, the tribonacci spectrum spans $[-1.569, 1.569]$ with a prominent
central gap of width $0.627$ and hierarchical sub-gaps, consistent with
Cantor-set spectral structure. The inverse participation ratio (IPR) of
mid-spectrum eigenstates is $\approx 0.087$, far above the extended-state
value $1/N = 0.002$, confirming multifractal critical behavior.

\section{The DNLS Model}

\subsection{Equation of Motion}

We study the discrete nonlinear Schr\"{o}dinger equation on the
$n$-bonacci substitution chain:
\begin{equation}
  i\,\frac{d\psi_j}{dt} = -\sum_{\langle j,k\rangle} t_{jk}\,\psi_k
    + \lambda\,|\psi_j|^2\,\psi_j,
  \label{eq:DNLS}
\end{equation}
where $t_{jk}$ are the hopping amplitudes dictated by the substitution
structure (Section~2.3), $\lambda > 0$ is the nonlinear coupling strength,
and the sum is over nearest neighbors $\langle j,k \rangle$ on the chain.
This is the standard focusing cubic DNLS, the discrete analog of the
Gross--Pitaevskii equation \cite{Eilbeck1985,Kevrekidis2009}.

\subsection{Amplitude Envelope Ansatz}

In the linear limit ($\lambda = 0$), the mid-gap eigenstate $\psi^{(0)}$
of the tribonacci Hamiltonian serves as our initial condition. For physical
motivation, we note that within the dm$^3$ framework \cite{Grossi2026_Vol1},
the natural amplitude envelope along the chain is
\[
  A_k \sim \eta^{-k}, \quad k \geq 0,
\]
where $\eta$ is the tribonacci constant. This is precisely the \emph{geometric
weight} $w_k = \eta^{-k}$ whose strict antitonicity (larger $k$ implies
smaller amplitude) is Lean~4-verified in \texttt{TribonacciMeasure.lean}
\cite{AXLE2026}:
\begin{equation}
  w_k = \eta^{-k}, \quad w_0 > w_1 > w_2 > \cdots > 0.
  \label{eq:weight}
\end{equation}
The mid-gap eigenstate of the tribonacci Hamiltonian inherits this decay
structure from the substitution geometry, making it the natural initial
condition for studying nonlinear perturbations of the geometric weight.

\subsection{Numerical Implementation}

Equation~\eqref{eq:DNLS} conserves the $\ell^2$ norm
$\mathcal{N} = \sum_j |\psi_j|^2$. We split $\psi_j = u_j + iv_j$ and
integrate the real $2N$-dimensional system with \texttt{scipy.integrate.solve\_ivp}
(RK45, relative tolerance $10^{-7}$, absolute tolerance $10^{-9}$, maximum
step $0.1$). Norm conservation was verified to hold to $< 10^{-6}$ throughout
all runs, confirming numerical accuracy.

\paragraph{Initial conditions.} For each chain, we compute the eigenstate
$\psi^{(0)}$ of the linear Hamiltonian with eigenvalue closest to $E = 0$
(the mid-gap state). This is the most multifractally critical state and
provides a physically meaningful, chain-specific initial condition.
For the tribonacci chain: mid-gap eigenvalue $= 0.000000$, IPR $= 0.0820$.
For the Fibonacci chain: mid-gap eigenvalue $= 0.000915$, IPR $= 0.0210$.

\paragraph{Localization measure.} We track the normalized IPR:
\[
  \mathrm{IPR}(t) = \frac{\sum_j |\psi_j(t)|^4}{\left(\sum_j |\psi_j(t)|^2\right)^2}.
\]
Extended states have IPR $\sim 1/N \to 0$; perfectly localized states have
IPR $= 1$; multifractal critical states lie in between.

\section{Results}

\subsection{Linear Baseline}

In the linear limit ($\lambda = 0$) the IPR of the mid-gap state is fixed by
the eigenstate structure:
\[
  \mathrm{IPR}_{\mathrm{trib}}(0) = 0.0820, \quad
  \mathrm{IPR}_{\mathrm{fib}}(0) = 0.0210.
\]
The ratio $\approx 3.91$ reflects the stronger multifractality of the
tribonacci chain. This is consistent with the larger fractal dimensions
computed in \cite{Krebbekx2023} for Rauzy substitution spectra: the tribonacci
chain's internal space (the Rauzy fractal) has higher dimension than the
Fibonacci chain's Cantor set, producing more strongly clustered eigenstates.

\subsection{Nonlinear Dynamics: IPR vs.\ $\lambda$}

Table~\ref{tab:ipr} and Figure~\ref{fig:ipr} (schematic; numerical values
exact) show the IPR at time $T = 50$ as a function of $\lambda$.

\begin{table}[h]
\centering
\caption{IPR at $T = 50$ for Fibonacci ($n=2$) and Tribonacci ($n=3$) chains,
$N = 500$ sites, $t = 0.5$. ``Ratio'' is IPR$_{\rm trib}$/IPR$_{\rm fib}$.}
\label{tab:ipr}
\begin{tabular}{r r r r}
\toprule
$\lambda$ & IPR$_{\rm trib}$ & IPR$_{\rm fib}$ & Ratio \\
\midrule
0.0  & 0.0820 & 0.0210 & 3.91 \\
0.5  & 0.0812 & 0.0168 & 4.82 \\
1.0  & 0.0789 & 0.0116 & 6.81 \\
1.5  & 0.0782 & 0.0090 & 8.65 \\
2.0  & 0.0704 & 0.0088 & 8.00 \\
3.0  & 0.0835 & 0.0115 & 7.28 \\
4.0  & 0.0587 & 0.0153 & 3.83 \\
5.0  & 0.0430 & 0.0173 & 2.48 \\
7.0  & 0.0316 & 0.0132 & 2.40 \\
10.0 & 0.0271 & 0.0082 & 3.30 \\
\bottomrule
\end{tabular}
\end{table}

\subsection{Differential Nonlinear Robustness}

The key finding is visible in the normalized IPR ratio
$\mathrm{IPR}(\lambda)/\mathrm{IPR}(0)$:

\begin{center}
\begin{tabular}{r r r}
\toprule
$\lambda$ & Tribonacci & Fibonacci \\
\midrule
0.0 & 1.000 & 1.000 \\
0.5 & 0.990 & 0.803 \\
1.0 & 0.962 & 0.552 \\
1.5 & 0.953 & 0.431 \\
2.0 & 0.859 & 0.420 \\
5.0 & 0.524 & 0.825 \\
10.0 & 0.330 & 0.392 \\
\bottomrule
\end{tabular}
\end{center}

At moderate coupling ($\lambda \leq 2$), the Fibonacci mid-gap state
delocalizes rapidly: IPR falls to $\sim 43\%$ of its linear value at
$\lambda = 1.5$. The tribonacci mid-gap state is nearly pinned: IPR falls
only to $\sim 95\%$ of its linear value at the same coupling. This is the
central result: \emph{the tribonacci chain's more strongly multifractal
critical states exhibit substantially greater resistance to nonlinear
delocalization in the moderate-coupling regime.}

At large coupling ($\lambda \gtrsim 5$), both chains delocalize, but the
tribonacci chain remains $\sim 2.4$--$3.3\times$ more localized throughout
the scanned range.

\begin{remark}[Non-Monotone Behavior at $\lambda \approx 3$]
The tribonacci IPR shows a local increase near $\lambda = 3$ (from $0.070$
at $\lambda = 2$ to $0.084$ at $\lambda = 3$). This non-monotone feature may
reflect a nonlinear resonance between the self-trapping threshold and the
energy of the mid-gap state. Understanding this feature requires longer
integration times and finite-size scaling analysis, which we leave to future
work.
\end{remark}

\subsection{Physical Interpretation}

The mechanism underlying differential nonlinear robustness is the following.
The tribonacci chain's stronger multifractality means its mid-gap eigenstate
has a more strongly hierarchical spatial structure: amplitude is concentrated
on a smaller effective number of sites (higher IPR). This spatial hierarchy
acts as an effective disorder landscape that the nonlinear term must overcome
to spread the wavepacket. The Fibonacci chain's less hierarchical eigenstate
is more susceptible to this spreading because its effective localization is
weaker to begin with.

More precisely: the nonlinear self-trapping transition in DNLS systems occurs
when the nonlinear energy $\lambda \langle |\psi|^4 \rangle = \lambda \cdot
\mathrm{IPR} \cdot \mathcal{N}^2$ is comparable to the bandwidth. The higher
initial IPR of the tribonacci state means it enters the self-trapping regime
at a lower effective $\lambda$---but because it is already strongly localized,
the self-trapping has less work to do. The Fibonacci state, being more
extended, requires more nonlinear energy to localize but also loses its initial
localization more easily under nonlinear perturbation.

\section{Connection to the dm$^3$ Framework}

The amplitude decay envelope $w_k = \eta^{-k}$ that appears in the dm$^3$
contact-geometric framework \cite{Grossi2026_Vol1,Grossi2026_Wavenumber6}
is identical to the Perron--Frobenius weight of the tribonacci companion
matrix. Within the dm$^3$ framework, this weight is not introduced by hand
but is forced by the three-axis decomposition of the contact 3-manifold
$(\mathcal{M}, \alpha)$ \cite{Grossi2026_Vol1,Grossi2026_Tribonacci}: two
axes in the contact distribution $\ker\alpha$ and one Reeb axis combine to
produce the minimal three-term recurrence, whose dominant eigenvalue is
precisely $\eta$.

The present paper does not require the full dm$^3$ framework; the substitution
chain and DNLS model stand independently. However, the formal verification of
$\eta > 1$ and $w_k \searrow 0$ in \texttt{TribonacciMeasure.lean}
\cite{AXLE2026} provides Lean~4-verified support for the mathematical
well-posedness of the amplitude envelope used as initial condition.

We note the following items whose proof in the dm$^3$ context remains open
and are not claimed here: (i) the connection between the Rauzy substitution
chain and the contact 3-manifold geometry; (ii) the `sorry'-carrying theorems
in \texttt{dm3\_Criticality\_gtct.lean}; (iii) the g$_{33}$ stability
threshold conjecture. These are tracked as open problems in AXLE
\cite{AXLE2026} (Issue \#13 and the sorry roadmap).

\section{Related Work and Gap Claim}

The closest existing work is \cite{Lahini2009}, who studied DNLS dynamics on
Fibonacci chains experimentally in photonic lattices and found subdiffusive
spreading governed by the multifractal exponent. Their finding---that stronger
multifractality leads to slower spreading---is qualitatively consistent with
our result, and provides an experimental framework into which tribonacci
measurements could in principle be inserted.

Ref.~\cite{Krebbekx2023} established the linear tribonacci tight-binding
model and its spectral properties but did not consider nonlinear dynamics.
To our knowledge, no prior work has studied DNLS or Gross--Pitaevskii
dynamics on a tribonacci or higher-$n$-bonacci substitution chain. The gap
claim is narrow and specific: nonlinear dynamics on tribonacci chains. We do
not claim novelty for the individual components (DNLS, tribonacci spectrum,
multifractality), only for their combination.

\section{Open Questions and Future Work}

The results presented here are a numerical first step. The following questions
are immediate:

\begin{enumerate}

\item \textbf{Longer time evolution.} Our integration runs to $T = 50$.
Subdiffusive spreading in quasiperiodic DNLS typically becomes visible on
timescales $T \sim 10^3$--$10^5$ \cite{Flach2010}. Whether the tribonacci
advantage persists at longer times is the most urgent open question.

\item \textbf{Finite-size scaling.} The IPR ratio $\sim 4$ between tribonacci
and Fibonacci mid-gap states should be checked at $N = 200, 500, 1000, 2000$
to confirm it reflects a bulk property and not finite-size artifact.

\item \textbf{Spreading exponent.} Computing the mean-square displacement
$\langle r^2 \rangle(t) \sim t^{\alpha}$ for both chains would allow
comparison with the Fibonacci spreading exponent established in
\cite{Lahini2009} and test whether the tribonacci exponent is genuinely
smaller.

\item \textbf{Self-trapping threshold.} A systematic scan over initial
amplitudes would determine the critical nonlinearity $\lambda_c^{(n)}$
below which spreading is subdiffusive and above which self-trapping occurs,
for $n = 2, 3$, and $4$.

\item \textbf{Lean~4 formalization.} The full DNLS model and the IPR
inequality $\mathrm{IPR}_{\rm trib}(0) > \mathrm{IPR}_{\rm fib}(0)$
(which follows from the spectral dimension inequality proved in
\cite{Krebbekx2023}) could in principle be added to the AXLE sorry roadmap
as formal verification targets.

\end{enumerate}

\section{Conclusion}

We have shown numerically that the tribonacci substitution chain ($n=3$)
exhibits substantially greater resistance to nonlinear delocalization than the
Fibonacci chain ($n=2$) when both are driven by a cubic DNLS nonlinearity
starting from mid-gap critical states. The tribonacci mid-gap state, with an
IPR $\approx 4\times$ larger than its Fibonacci counterpart in the linear
limit, retains $> 95\%$ of its localization at nonlinear coupling $\lambda =
1.5$ where the Fibonacci state has lost $\sim 57\%$ of its localization. We
term this \emph{differential nonlinear robustness} and identify its mechanism
as the stronger effective spatial hierarchy of the tribonacci multifractal
eigenstate.

The tribonacci amplitude decay constant $\eta \approx 1.839287$ is formally
verified to satisfy $\eta > 1$ and to produce a strictly antitone weight
sequence $\{\eta^{-k}\}$ in Lean~4/Mathlib4 \cite{AXLE2026}, providing
machine-checked support for the mathematical well-posedness of the amplitude
envelope.

These results constitute, to our knowledge, the first numerical study of
DNLS dynamics on a tribonacci substitution chain and identify a concrete,
measurable difference between $n=2$ and $n=3$ quasicrystal models under
nonlinear perturbation.

\section*{Acknowledgements}

The author thanks the Lean~4/Mathlib4 community for the formal verification
infrastructure used in \texttt{TribonacciMeasure.lean}. Numerical computations
used Python~3 with NumPy, SciPy, and standard scientific libraries.

\begin{thebibliography}{99}

\bibitem{Kohmoto1983}
M.~Kohmoto, L.~P.~Kadanoff, and C.~Tang,
``Localization Problem in One Dimension: Mapping and Escape,''
\textit{Phys.\ Rev.\ Lett.} \textbf{50}, 1870 (1983).

\bibitem{Ostlund1983}
S.~\"{O}stlund, R.~Pandit, D.~Rand, H.~J.~Schellnhuber, and E.~D.~Siggia,
``One-Dimensional Schr\"{o}dinger Equation with an Almost Periodic Potential,''
\textit{Phys.\ Rev.\ Lett.} \textbf{50}, 1873 (1983).

\bibitem{Luck1989}
J.~M.~Luck,
``Cantor spectra and scaling of gap widths in deterministic aperiodic systems,''
\textit{Phys.\ Rev.\ B} \textbf{39}, 5834 (1989).

\bibitem{Mace2017}
N.~Mac\'{e}, A.~Jagannathan, and F.~Pi\'{e}chon,
``Fractal dimensions of wave functions and local spectral measures
on the Fibonacci chain,''
\textit{Phys.\ Rev.\ B} \textbf{93}, 205134 (2016).

\bibitem{Kraus2012}
Y.~E.~Kraus, Y.~Lahini, Z.~Ringel, M.~Verbin, and O.~Zilberberg,
``Topological States and Adiabatic Pumping in Quasicrystals,''
\textit{Phys.\ Rev.\ Lett.} \textbf{109}, 106402 (2012).

\bibitem{Lang2012}
L.-J.~Lang, X.~Cai, and S.~Chen,
``Edge States and Topological Phases in One-Dimensional Optical Superlattices,''
\textit{Phys.\ Rev.\ Lett.} \textbf{108}, 220401 (2012).

\bibitem{Krebbekx2023}
N.~J.~Krebbekx, A.~Moustaj, K.~Dajani, and C.~Morais~Smith,
``Multifractal properties of tribonacci chains,''
\textit{Phys.\ Rev.\ B} \textbf{108}, 104204 (2023).

\bibitem{Bellissard1989}
J.~Bellissard, B.~Iochum, E.~Scoppola, and D.~Testard,
``Spectral properties of one dimensional quasi-crystals,''
\textit{Commun.\ Math.\ Phys.} \textbf{125}, 527 (1989).

\bibitem{Damanik2008}
D.~Damanik and A.~Gorodetski,
``Hyperbolicity of the trace map for the weakly coupled Fibonacci Hamiltonian,''
\textit{Nonlinearity} \textbf{22}, 123 (2009).

\bibitem{Rauzy1982}
G.~Rauzy,
``Nombres alg\'{e}briques et substitutions,''
\textit{Bull.\ Soc.\ Math.\ France} \textbf{110}, 147 (1982).

\bibitem{Lahini2009}
Y.~Lahini, R.~Pugatch, F.~Pozzi, M.~Sorel, R.~Morandotti, N.~Davidson,
and Y.~Silberberg,
``Observation of a Localization Transition in Quasiperiodic Photonic Lattices,''
\textit{Phys.\ Rev.\ Lett.} \textbf{103}, 013901 (2009).

\bibitem{Jurgensen2021}
M.~J\"{u}rgensen, S.~Mukherjee, and M.~C.~Rechtsman,
``Quantized nonlinear Thouless pumping,''
\textit{Nature} \textbf{596}, 63 (2021).

\bibitem{Verbin2015}
M.~Verbin, O.~Zilberberg, Y.~Lahini, Y.~E.~Kraus, and Y.~Silberberg,
``Topological pumping over a photonic Fibonacci quasicrystal,''
\textit{Phys.\ Rev.\ B} \textbf{91}, 064201 (2015).

\bibitem{Eilbeck1985}
J.~C.~Eilbeck, P.~S.~Lomdahl, and A.~C.~Scott,
``The discrete self-trapping equation,''
\textit{Physica D} \textbf{16}, 318 (1985).

\bibitem{Kevrekidis2009}
P.~G.~Kevrekidis,
\textit{The Discrete Nonlinear Schr\"{o}dinger Equation}.
Springer, Berlin (2009).

\bibitem{Flach2010}
S.~Flach, M.~V.~Ivanchenko, and O.~I.~Kanakov,
``Spreading of wave packets in one-dimensional disordered lattices,''
\textit{Phys.\ Rev.\ E} \textbf{82}, 036219 (2010).

\bibitem{Grossi2026_DNLS}
P.~Nogueira~Grossi,
``Differential Nonlinear Robustness of Critical States in Fibonacci and
Tribonacci Substitution Chains.''
G6~LLC. Zenodo: \href{https://doi.org/10.5281/zenodo.20026943}{10.5281/zenodo.20026943} (2026).

\bibitem{Grossi2026_Vol1}
P.~Nogueira~Grossi,
\textit{Principia Orthogona, Volume~I: The Mathematics of Generative Transitions}.
G6~LLC. Zenodo: \href{https://doi.org/10.5281/zenodo.19117400}{10.5281/zenodo.19117400} (2026).

\bibitem{Grossi2026_Wavenumber6}
P.~Nogueira~Grossi,
``Wavenumber 6: The Orthogenetic Stability Generator of Nested Infinities.''
G6~LLC. Zenodo: \href{https://doi.org/10.5281/zenodo.19199474}{10.5281/zenodo.19199474} (2026).

\bibitem{Grossi2026_Tribonacci}
P.~Nogueira~Grossi,
``Convergent Derivations of the Tribonacci Coherence Threshold:
Contact Geometry, Rational Quantum Mechanics, and the SL(3,Z) Algebra.''
G6~LLC (2026). Manuscript.

\bibitem{AXLE2026}
P.~Nogueira~Grossi,
\textit{AXLE: Lean~4 Formal Verification Engine}.
GitHub: \href{https://github.com/TOTOGT/AXLE}{github.com/TOTOGT/AXLE} (2026).
[\texttt{TribonacciMeasure.lean}: $\eta > 1$ and strict antitonicity of
$\eta^{-k}$, proved without \texttt{sorry}.]

\end{thebibliography}

\end{document}
